\section{Pricing from one cumulative ledger}
\label{sec:ledger}

Once the transport $x$ is known, pricing does not require a new integration
for every strike.  One cumulative tail integral contains nearly every
quantity required by a trader or a fitter.

\subsection{The upper-share ledger}

\begin{definition}[Upper-share ledger]
\label{def:ledger}
For an admissible slice define
\begin{equation}
G(z)=\int_z^\infty e^{x(t)}\rho(t)\,\dd t
=\E[Y\1\{Z>z\}].
\label{eq:ledger}
\end{equation}
\end{definition}

The ordinary upper-tail probability is $1-\Lambda(z)$.  The ledger $G(z)$
weights the same outcomes by the terminal gross return $Y=e^X$.  It is the
fraction of the unit forward contributed by scenarios above rank
$\Lambda(z)$.  It satisfies
\begin{equation}
G(-\infty)=1,\qquad G(+\infty)=0,
\qquad G'(z)=-e^{x(z)}\rho(z)<0.
\label{eq:ledgerprops}
\end{equation}
For a strike $k$, let $z_k$ be the unique point at which the transported
log-return reaches the strike, $x(z_k)=k$, and set
$u_k=\Lambda(z_k)=F_X(k)$.  Then $1-u_k$ is the exercise probability.

\begin{proposition}[Call, put, digital, and density]
\label{prop:pricing}
For every admissible slice,
\begin{align}
c(k)&=G(z_k)-e^k(1-u_k),
\label{eq:call}\\
P(k)&=e^ku_k-(1-G(z_k)),
\label{eq:put}\\
c'(k)&=-e^k(1-u_k),
\label{eq:cprime}\\
c''(k)&=-e^k(1-u_k)+e^k\frac{\rho(z_k)}{x'(z_k)},
\label{eq:csecond}\\
f_X(k)&=e^{-k}\{c''(k)-c'(k)\}
=\frac{\rho(z_k)}{x'(z_k)}.
\label{eq:bl_division}
\end{align}
\end{proposition}

\begin{proof}
The call is exercised on $Z>z_k$.  Splitting its payoff into the received
asset and the paid strike gives
\[
c(k)=\E[Y\1\{Z>z_k\}]-e^k\Prob(Z>z_k)
=G(z_k)-e^k(1-u_k).
\]
The put follows on the complementary event.  When the call is differentiated,
the moving-root term is
\[
\{G'(z_k)+e^k\rho(z_k)\}\frac{\dd z_k}{\dd k}=0
\]
because $G'(z_k)=-e^{x(z_k)}\rho(z_k)=-e^k\rho(z_k)$.  This gives
\eqref{eq:cprime}.  Since $\dd z_k/\dd k=1/x'(z_k)$, one more derivative
gives \eqref{eq:csecond} and \eqref{eq:bl_division}.
\end{proof}

The pricing rule is worth retaining in words:
\begin{equation}
\boxed{\quad\text{call value}=\text{upper asset share}
-\text{strike}\times\text{exercise probability}.\quad}
\label{eq:pricingwords}
\end{equation}
The digital is $1-u_k$.  The density is one division of two quantities
already present in the transport.  Neither is obtained by differencing noisy
option prices.

More generally, any integrable terminal payoff is a rank integral:
\begin{equation}
\E[H(Y)]=\int_{-\infty}^{\infty}H(e^{x(z)})\rho(z)\,\dd z.
\label{eq:generalpayoff}
\end{equation}
In particular, the log-contract variance is
\begin{equation}
w_{\rm vs}=-2\E[\log Y]
=-2\int_{-\infty}^{\infty}x(z)\rho(z)\,\dd z.
\label{eq:varswap}
\end{equation}
This is the distribution-side form of the usual option-strip replication
\citep{Neuberger1994,DemeterfiDermanKamalZou1999}.  Under continuous paths and
continuous monitoring it is the fair variance-swap strike.  Jumps and discrete
sampling require their standard corrections.

\subsection{A complete option calculation}

In the constant-speed reference slice, choose $s=\MacToyScale$, which gives a
six-month ATM volatility of twenty percent after exact repricing.  The shift
is $m=\MacToyMu$.  For $k=0.10$, the strike rank is
$u_k=\MacToyOtmPercentile\%$.  The ledger and cash legs are
\[
G(z_k)=\MacToyShareOtm,
\qquad e^k(1-u_k)=\MacToyCashOtm.
\]
The normalized call is therefore
\[
c(0.10)=\MacToyShareOtm-\MacToyCashOtm=\MacToyCallOtm,
\]
which corresponds to \MacToyIvOtmPct\% implied volatility.  Every market
option uses the same three operations: locate the strike rank, read the
ledger, and subtract the strike cash leg.

\subsection{Reading level, skew, and curvature at the forward}

This subsection is not needed for pricing.  It explains exactly which local
features of the fitted law generate the three familiar at-the-money (ATM)
smile statistics.  Let $z_*$ solve $x(z_*)=0$ and define
\begin{equation}
c_0=G(z_*)-(1-u_*),\qquad
u_*=\Lambda(z_*),\qquad
f_*=\frac{\rho(z_*)}{x'(z_*)}.
\label{eq:atm_inputs}
\end{equation}
The ATM Black variance and half-standard-deviation are
\[
w_0=4\left[\PhiN^{-1}\left(\frac{1+c_0}{2}\right)\right]^2,
\qquad d=\frac{\sqrt{w_0}}2.
\]
Define the digital gap
\begin{equation}
\Delta=u_*-\PhiN(d).
\label{eq:digitalgap}
\end{equation}

\begin{proposition}[Exact ATM handles]
\label{prop:atm}
The total-variance smile $\Black(k,w(k))=c(k)$ satisfies
\begin{align}
w(0)&=w_0,
&w'(0)&=\frac{2\sqrt{w_0}}{\phin(d)}\Delta,
\label{eq:atm_w}\\
w''(0)&=\frac{2\sqrt{w_0}}{\phin(d)}f_*-2
+\left(\frac1{2w_0}+\frac18\right)w'(0)^2.
\label{eq:atm_wpp}
\end{align}
In volatility units, $\sigma(k)=\sqrt{w(k)/\tau}$,
\begin{align}
\sigma(0)&=\sqrt{w_0/\tau},
&\sigma'(0)&=\frac{\Delta}{\phin(d)\sqrt\tau},
\label{eq:atm_sigma1}\\
\sigma''(0)&=\frac1{\sqrt\tau}
\left[\frac{f_*}{\phin(d)}-\frac1{\sqrt{w_0}}
+\frac{\sqrt{w_0}}4\left(\frac{\Delta}{\phin(d)}\right)^2\right].
\label{eq:atm_sigma}
\end{align}
\end{proposition}

The derivation is in \cref{app:atmproof}.  The interpretation is simpler than
the formulas:
\begin{itemize}
\item ATM option value fixes the volatility level;
\item the digital gap $\Delta$ fixes the skew;
\item density at the forward supplies the additional information needed for
curvature.
\end{itemize}
For a matched Black law, $u_*=\PhiN(d)$ and the skew is zero.  A negative
equity skew means $u_*<\PhiN(d)$: the model assigns more probability above
the forward than the flat-Black law with the same ATM call value.

\subsection{The analytic calibration Jacobian}
\label{sec:jacobian}

The call depends on a parameter in two ways: the ledger changes, and the rank
at which the strike is reached moves.  These effects appear to require a root
sensitivity for every strike.  In fact the moving-root term cancels exactly.

\begin{theorem}[Envelope cancellation]
\label{thm:envelope}
For every coordinate $\theta_j$,
\begin{equation}
\frac{\partial c(k;\theta)}{\partial\theta_j}
=\left.\frac{\partial G(z;\theta)}{\partial\theta_j}
\right|_{z=z_k(\theta)}.
\label{eq:envelope}
\end{equation}
\end{theorem}

\begin{proof}
Differentiate \eqref{eq:call}.  The coefficient of
$\partial z_k/\partial\theta_j$ is
$G'(z_k)+e^k\rho(z_k)$, which is zero by \eqref{eq:ledgerprops} and
$x(z_k)=k$.
\end{proof}

Let $b_j(u)=\partial g(u)/\partial\theta_j$ and
$\bar x=x-m$.  For fixed tail exponents,
\begin{align}
\frac{\partial\bar x(z)}{\partial\theta_j}
&=\int_0^z x'(t)b_j(\Lambda(t))\,\dd t,
\label{eq:pass1}\\
\frac{\partial m}{\partial\theta_j}
&=-\int_{-\infty}^{\infty}e^{x(z)}
\frac{\partial\bar x(z)}{\partial\theta_j}\rho(z)\,\dd z,
\label{eq:pass2}\\
\frac{\partial G(z)}{\partial\theta_j}
&=\int_z^\infty e^{x(t)}
\left(\frac{\partial m}{\partial\theta_j}
+\frac{\partial\bar x(t)}{\partial\theta_j}\right)
\rho(t)\,\dd t.
\label{eq:pass3}
\end{align}
These are three cumulative passes: differentiate transport, differentiate the
normalizing shift, then integrate the ledger sensitivity backwards.  The
denominator in the derivative of $m$ is one because the accepted slice is
already martingale-normalized.  In raw coordinates,
\[
b_L(u)=1-u,\qquad b_R(u)=u,\qquad
b_{a_n}(u)=P_n(1-2u).
\]
The same arrays supply the variance-swap sensitivity:
\begin{equation}
\frac{\partial w_{\rm vs}}{\partial\theta_j}
=-2\int_{-\infty}^{\infty}
\left(\frac{\partial m}{\partial\theta_j}
+\frac{\partial\bar x(z)}{\partial\theta_j}\right)
\rho(z)\,\dd z.
\label{eq:varswaprow}
\end{equation}

If a tail exponent is estimated, it adds the columns
\begin{equation}
b_{\alpha_-}(u)=-\log\ell_-(u),
\qquad b_{\alpha_+}(u)=-\log\ell_+(u)
\label{eq:alphacolumns}
\end{equation}
to the derivative of $\log x'$.  These columns grow in the remote ranks and
are weakly identified by a finite strike strip.  The reference policy fixes
$\alpha_\pm$ and reports scenario fits instead.

\begin{figure}[htbp]
\centering
\paperfig{0.92\textwidth}{fig_jacobian.pdf}
\caption{Analytic price sensitivities for the fitted SPY December 2026 slice.
Left: one normalized sensitivity curve per parameter.  Right: disagreement
with an independent central difference.  The comparison includes transport,
normalization, ledger integration, root finding, and interpolation.}
\label{fig:jacobian}
\end{figure}

For $P=N+1$ body and endpoint parameters and $n_{\rm grid}$ nodes, the entire
sensitivity pass costs $O(Pn_{\rm grid})$.  Its cost is essentially independent
of the number of quoted strikes.  This is the practical value of the envelope
cancellation: no slice rebuild and no repeated strike-root solve is needed
for each finite-difference column.
